\chapter{Fiabilisation de la plateforme Protova1 : mécanique et odométrie}

\section{Contexte}
% TODO : Protova1, première génération : base mécanique vieillissante,
% odométrie inexploitable à cause des rebonds de l'encodeur.

\section{Problématique}
% TODO : comment redonner à Protova1 une base saine (mécanique + mesure de
% déplacement fiable) pour en faire à terme un second véhicule téléopérable ?

\section{Mission 1 : Conception et impression 3D de la nouvelle base mécanique}
\objectif
% TODO : remplacer la base existante ; contraintes d'intégration (fixations
% Jetson/lidar/batterie, rigidité, masse).
\methodes
% TODO : CAO (mêmes réflexes que tes conceptions FORVIA : cahier des charges,
% itérations de solutions, choix justifié), impression 3D : matériau,
% orientation, tolérances, assemblage.
\resultats
% TODO : photos avant/après, masse, qualité d'assemblage.
\discussion
% TODO : itérations ratées/réussies, leçons (tolérances d'impression...).

\section{Mission 2 : Chaîne odométrique --- filtrage de l'encodeur}
\objectif
% TODO : signaux du KY-040 inexploitables (rebonds) -> compter proprement les
% ticks pour l'odométrie.
\methodes
% TODO : analyse du signal brut, conception du conditionnement : filtre RC +
% trigger de Schmitt 74HC14, schéma électronique, dimensionnement justifié
% (choix R, C, seuils d'hystérésis), banc de test Arduino en quadrature.
\resultats
% TODO : LE chiffre : 1 tour = 80 ticks, stable et reproductible — mets
% l'oscillogramme avant/après filtrage, c'est très parlant.
\discussion
% TODO : limites (résolution, vitesse max), alternative logicielle écartée.

\section{Vers la téléopération de Protova1}
% TODO : perspective — ce qui reste à faire pour rendre Protova1 téléopérable
% comme Protova2 (annonce des perspectives de la conclusion générale).

\section{Conclusion}
% TODO : 4-5 lignes de bilan du chapitre.
